\chapter{Cross-Region Crypto Networks}\label{ch:nw:cross-region-crypto-networks}

Crypto liquidity is spread over a few cloud regions, and a firm that trades across them needs its messages to cross oceans. In April 2026
a round trip from Tokyo to Frankfurt over the cloud provider's own backbone, by VPC peering, took a median \qty{224.3}{\milli\second};
over the same provider's Direct Connect joined to a specialised network, \qty{136.5}{\milli\second}. Light in straight glass along the
geodesic would need \qty{91.3}{\milli\second}. The provider's backbone is built for ``resilience and scale'', its blog says, and it takes the
long way; the specialist takes a shorter one; and the firm that cannot choose sends the same message down both and keeps the first to arrive.

Chapters~16 to~18 worked inside one region. This chapter connects the regions: the hubs and their distances, the backbone and the internet
against private lines and specialised providers, and \omterm{def:nw:cross-region-crypto-networks:idem}{route racing}, with the duplicate orders it creates and what a venue's rules do with
them. The published figures are cited and dated; the racing is a labelled simulation.

\section{The hubs: Tokyo, Singapore, Hong Kong, London, Virginia}

\begin{definition}[Path inflation]\label{def:nw:cross-region-crypto-networks:infl}
The \emph{path inflation}\index{path inflation} of a route between two places is its measured round trip divided by the round trip light would
need along the geodesic in the route's medium (for long-haul networks, fibre): chapter~10's \omterm{def:nw:the-north-american-map:factor}{route factor}, for round trips over networks whose
medium and path are not published.
\end{definition}

The hubs are chapter~17's: Tokyo, where AWS's ap-northeast-1 region hosts Binance and BitMEX and Hyperliquid's validators; Singapore, Bybit's
region; Hong Kong; London and Frankfurt in Europe; Northern Virginia in the United States. The provider publishes only the regions' names and
their numbers of zones (four in Tokyo, six in Northern Virginia, three in the others here), so the chapter places each at its city, and
Northern Virginia at Ashburn, an assumption flagged in the site table.

\begin{table}[ht]
\centering\small
\begin{tabular}{lrrrrrrr}
\toprule
From Tokyo to & Geodesic & Vacuum & Fibre & Backbone & Specialist & \multicolumn{2}{c}{\omterm{def:nw:cross-region-crypto-networks:infl}{Path inflation}} \\
 & (km) & RTT (ms) & RTT (ms) & (ms) & (ms) & backbone & specialist \\
\midrule
Hong Kong & 2\,883 & 19.2 & 28.1 & -- & -- & -- & -- \\
Singapore & 5\,311 & 35.4 & 51.8 & 68.0 & 65.4 & 1.31 & 1.26 \\
N. Virginia & 10\,895 & 72.7 & 106.3 & 150.2 & 135.4 & 1.41 & 1.27 \\
Frankfurt & 9\,359 & 62.4 & 91.3 & 224.3 & 136.5 & 2.46 & 1.50 \\
London & 9\,586 & 64.0 & 93.5 & 210.0 & 139.6 & 2.25 & 1.49 \\
Stockholm & 8\,194 & 54.7 & 79.9 & 245.3 & 124.3 & 3.07 & 1.56 \\
\bottomrule
\end{tabular}
\caption{Tokyo's hubs: \omterm{def:nw:the-north-american-map:geodesic}{geodesic distances} between the region cities and the round trips light needs along them, against the median round trips AWS
published for April 2026 over its backbone (VPC peering) and over Direct Connect with a specialised network, and the \omterm{def:nw:cross-region-crypto-networks:infl}{path inflation} of each against
the fibre floor. Data: \texttt{nw\_race.hub\_rows()}.}\label{tab:nw:cross-region-crypto-networks:hubs}
\end{table}

The table's pattern is the chapter's subject. To Singapore and Virginia both networks run 26 to 41\% above straight glass. To Europe the
backbone runs 2.2 to 3.1 times the floor, the specialist 1.5 times. Neither publishes its path; the inflation says the backbone goes a long
way round and the specialist a shorter one. The best European route from Tokyo, \qty{124.3}{\milli\second} to Stockholm, is still
more than twice the vacuum floor.

\begin{omfigure}
\begin{tikzpicture}
\begin{axis}[ybar,width=11.5cm,height=6cm,bar width=8pt,ymin=0,ymax=270,xtick=data,
  symbolic x coords={Singapore,N. Virginia,Frankfurt,London,Stockholm},ylabel={round trip from Tokyo (ms)},
  tick label style={font=\footnotesize},label style={font=\small},grid=major,grid style={gray!25},legend columns=3,
  legend style={font=\footnotesize,at={(0.5,-0.18)},anchor=north,draw=none},table/col sep=comma,enlarge x limits=0.15]
\addplot[fill=gray!40,draw=gray] table[x=hub,y=fibre_floor]{figdata/networks/19-cross-region-crypto-networks/hubs.csv};
\addplot[fill=omDef!55,draw=omDef] table[x=hub,y=specialist]{figdata/networks/19-cross-region-crypto-networks/hubs.csv};
\addplot[fill=omThm!45,draw=omThm] table[x=hub,y=backbone]{figdata/networks/19-cross-region-crypto-networks/hubs.csv};
\legend{fibre floor (geodesic),specialist network,provider backbone}
\end{axis}
\end{tikzpicture}

\omcaption{Round trips from Tokyo: the fibre floor along the geodesic, and the published medians over a specialised network and over the provider's
backbone (April 2026). Data: \texttt{fig\_race.py}, \cref{tab:nw:cross-region-crypto-networks:hubs}.}\label{fig:nw:cross-region-crypto-networks:hubs}
\end{omfigure}

\section{Cloud backbone, public internet and private lines}

\begin{definition}[Border Gateway Protocol, cloud backbone]\label{def:nw:cross-region-crypto-networks:bgp}
The \emph{Border Gateway Protocol}\index{Border Gateway Protocol} (BGP) is the internet's routing protocol between \omterm{def:nw:where-the-crypto-venues-live:as}{autonomous systems}: each exchanges
reachability information, including the list of \omterm{def:nw:where-the-crypto-venues-live:as}{autonomous systems} a route traverses, and chooses among routes by policy, not by latency. A
\emph{cloud backbone}\index{cloud backbone} is a provider's private network between its regions, which carries traffic between them (peering,
replication, its own services) under the provider's routing.
\end{definition}

Three kinds of path join two regions. The public internet, where each hop's operator chooses the next by business policy and BGP carries
reachability, not delay: short paths are a side effect. The provider's backbone, engineered for capacity and resilience, and in AWS's own blog not
the lowest latency. And private lines: \omterm{def:nw:long-haul-fibre:dark}{dark fibre} or wavelengths (chapter~13) and radio (chapter~14), bought directly or from a network that
specialises in low latency and connects to the provider's regions through its interconnection product, as the blog's specialist does.

\begin{dated}{2026-09}{Cross-region routes and duplicate rules, from public sources}\label{dat:nw:cross-region-crypto-networks:routes}
\begin{itemize}
\item AWS (June 2026): median TCP round trips from Tokyo on 28 April 2026, over VPC peering and over Direct Connect with Avelacom: Singapore 68.0 and
  65.4~ms; Virginia 150.2 and 135.4; Frankfurt 224.3 and 136.5; London 210.0 and 139.6; Stockholm 245.3 and 124.3.
\item Binance spot API: a client order identifier is unique among open orders; an order reusing one can be accepted only once the earlier order is
  filled, otherwise it is rejected (``Duplicate order sent.'').
\end{itemize}
\end{dated}

\section{Specialised low-latency providers and virtual-server hosts}

Between Tokyo and Europe the specialist of the blog saves 70 to 121~milliseconds of round trip over the backbone. For a strategy that trades a
Tokyo venue against a European one, that is the difference between seeing the other market's prices a tenth of a second late and seeing them in
time to matter. The specialist is a \omterm{def:nw:buying-connectivity:extranet}{financial extranet} (chapter~15) for cloud regions: private lines between the hubs, delivered into the
provider's regions through its interconnection product.

The \omterm{def:nw:where-the-crypto-venues-live:lottery}{instance lottery} of chapter~17 applies at both ends: the fastest cross-region route is wasted on a machine that sits far from the venue in
its own region.

\section{Racing the same message down several routes}

\begin{definition}[Route racing, idempotency key]\label{def:nw:cross-region-crypto-networks:idem}
\emph{Route racing}\index{route racing} sends the same message on several routes at once and uses the first copy to arrive. An \emph{idempotency
key}\index{idempotency key} is an identifier the sender attaches to every copy of one message (for orders, the client order identifier) so that the
receiver can recognise the copies and act on only one.
\end{definition}

\begin{proposition}[What racing buys]\label{prop:nw:cross-region-crypto-networks:race}
The first of $k$ copies arrives at the minimum of the routes' latencies. If the routes' delays were perfectly correlated, the minimum would be the
fastest route's delay and racing would buy nothing; the less correlated they are, the more the minimum's tail falls below every single route's
tail, since the minimum exceeds a level only when every route does.
\end{proposition}

\begin{proof}
$P(\min_i X_i > t) = P(X_1 > t, \dots, X_k > t)$, which equals the smallest single survival probability when the routes move together and falls to
their product when they are independent.
\end{proof}

The model races three one-way routes from Tokyo to Singapore: the specialist and the backbone at half their published round trips (32.7 and
\qty{34.0}{\milli\second}), each with an assumed 99th percentile 15\% above its median, and an internet route at \qty{35}{\milli\second} with a
heavier tail (an assumption). A Gaussian copula ties them with correlation $\rho$, a common cause such as a shared cable or congestion.

\begin{omfigure}
\begin{tikzpicture}
\begin{axis}[width=11cm,height=5.8cm,xmin=0,xmax=0.9,ymin=0,ymax=2300,xlabel={correlation of the routes' delays ($\rho$)},
  ylabel={p99 saved (\si{\micro\second})},tick label style={font=\footnotesize,/pgf/number format/1000 sep={\,}},label style={font=\small},grid=major,grid style={gray!25},
  legend pos=north east,legend style={font=\footnotesize},table/col sep=comma]
\addplot[omDef,thick,mark=*,mark size=1.4pt] table[x=rho,y=two]{figdata/networks/19-cross-region-crypto-networks/race.csv};
\addplot[omThm,thick,mark=square*,mark size=1.4pt] table[x=rho,y=three]{figdata/networks/19-cross-region-crypto-networks/race.csv};
\legend{two routes (specialist and backbone),three routes (with the internet)}
\end{axis}
\end{tikzpicture}

\omcaption{Simulation: the 99th-percentile one-way latency saved by racing, Tokyo to Singapore, against the best single route, as the routes'
correlation rises. At $\rho = 0$ two routes save 1.6 and three \qty{2.1}{\milli\second}; at $\rho = 0.9$ both save about
\qty{50}{\micro\second}. Data: \texttt{fig\_race.py}, \texttt{nw\_race.race\_rows()}.}\label{fig:nw:cross-region-crypto-networks:race}
\end{omfigure}

\omcode{code/firm/routerace/firm_routerace.py}{35}{50}{Correlated routes through a Gaussian copula, the first arrival, and the percentile saved
against the best single route.}\label{lst:nw:cross-region-crypto-networks:race}

Racing pays in the tail, and only while the routes fail independently. With independent routes the second route saves \qty{1623}{\micro\second} at
the 99th percentile and the third another \qty{518}{\micro\second}; at $\rho = 0.5$, 771 and 90; at 0.9, 53 and nothing. At an assumed 10~USD a
month per microsecond of 99th percentile, the third route is worth 5\,182~USD a month if the routes are independent, 902 at $\rho = 0.5$ and
nothing at 0.9: the price above which racing it does not pay.

What arrives second is a problem. The venue receives the same order several times; the \omterm{def:nw:cross-region-crypto-networks:idem}{idempotency key} must make it act once. Binance's spot
interface is typical of the rule to check: a client order identifier must be unique among open orders, and a copy is rejected as a duplicate only
while the first order is open. Once the first copy has filled, a late copy with the same identifier is a new order.

\begin{omfigure}
\begin{tikzpicture}
\begin{axis}[width=11cm,height=5.4cm,xmin=0,xmax=8000,ymin=0,ymax=2.1,xlabel={time for the first copy to fill (\si{\micro\second})},x tick label style={/pgf/number format/1000 sep={\,}},
  ylabel={duplicates accepted per order},tick label style={font=\footnotesize},label style={font=\small},grid=major,grid style={gray!25},
  table/col sep=comma,scaled x ticks=false]
\addplot[omMeth,thick,mark=*,mark size=1.4pt] table[x=fill_us,y=accepted]{figdata/networks/19-cross-region-crypto-networks/dups.csv};
\end{axis}
\end{tikzpicture}

\omcaption{Simulation: with three raced copies ($\rho = 0.5$) and a venue that rejects a duplicate identifier only while the first order is open, the
expected number of late copies accepted as new orders, against the time the first copy takes to fill. An order that fills on arrival is executed
three times. Data: \texttt{fig\_race.py}, \texttt{nw\_race.duplicate\_rows()}.}\label{fig:nw:cross-region-crypto-networks:dups}
\end{omfigure}

\begin{proposition}[Racing orders that fill]\label{prop:nw:cross-region-crypto-networks:dup}
Under a rule that rejects a repeated identifier only while the first order is open, a raced order that fills within time $f$ of the first copy's
arrival is executed again by every later copy that arrives after $f$. For orders that fill on arrival ($f = 0$), $k$ raced copies are $k$ orders.
\end{proposition}

\begin{proof}
A late copy that finds the first order filled finds no open order with its identifier and is accepted. With $f = 0$ every later copy is late.
\end{proof}

The model makes the danger concrete: three copies of an order that fills on arrival are three executions; one that takes a millisecond to fill
still produces 1.57 unwanted executions on average; eight milliseconds, 0.07. Racing is therefore safe for messages that are idempotent at the
venue (a cancel, a query, market data) or under a rule that remembers identifiers after the order has ended, and dangerous for taking orders under an
open-only rule. One Quant Book~13's idempotent processing is what the venue must do; this chapter's lesson is to read the venue's rule before racing.

\begin{omfigure}
\begin{tikzpicture}[font=\footnotesize,b/.style={draw,rounded corners=2pt,minimum height=0.55cm,align=center,inner sep=3pt},>=Stealth]
\node[b,fill=omDef!15] (s) at (0,0) {firm (Tokyo)\\ order, key K};
\node[b,fill=omThm!15,minimum height=1.6cm] (v) at (9.6,0) {venue\\ (Singapore)};
\draw[->,thick] (1.2,0.3) -- node[above,font=\scriptsize]{specialist: first, accepted} (8.7,0.5);
\draw[->,thick,gray] (1.2,0) -- node[above,font=\scriptsize]{backbone: second} (8.7,0);
\draw[->,thick,gray,dashed] (1.2,-0.3) -- node[below,font=\scriptsize]{internet: third} (8.7,-0.5);
\node[font=\scriptsize,align=left,anchor=west] at (10.6,0.3) {K open: reject copies};
\node[font=\scriptsize,align=left,anchor=west,omMeth] at (10.6,-0.3) {K filled: accept a new order};
\end{tikzpicture}

\omcaption{\omterm{def:nw:cross-region-crypto-networks:idem}{Route racing}, schematically: three copies of one order with one \omterm{def:nw:cross-region-crypto-networks:idem}{idempotency key}; the venue keeps the first and must recognise the others.
Under an open-only rule it rejects them while the order is open and accepts them as new orders once it has filled.}\label{fig:nw:cross-region-crypto-networks:scheme}
\end{omfigure}

\begin{method}[Racing routes between regions]\label{met:nw:cross-region-crypto-networks:race}
\begin{enumerate}
\item Measure each candidate route's latency distribution, not only its median, and the correlation of their delays (do they share a cable, a
  landing, a peering point?).
\item Race only routes whose tails are not too correlated, and price each extra route against the tail it saves.
\item Read the venue's duplicate rule. Race cancels, amendments and queries freely; race taking orders only if the venue remembers identifiers after
  an order ends, or if the firm's own logic tolerates a duplicate.
\item Log every copy's arrival (from the venue's acknowledgements) and re-measure the routes' correlation; it changes when a cable is cut or
  repaired.
\end{enumerate}
\end{method}

\section{Tutorial: hubs and races}\label{tut:nw:cross-region-crypto-networks:tut}

\begin{tutorial}
\textbf{Goal.} Measure \omterm{def:nw:cross-region-crypto-networks:infl}{path inflation} between the hubs and simulate \omterm{def:nw:cross-region-crypto-networks:idem}{route racing} with its duplicates.
\textbf{End state:} \cref{fig:nw:cross-region-crypto-networks:hubs,fig:nw:cross-region-crypto-networks:race,fig:nw:cross-region-crypto-networks:dups}
and \cref{tab:nw:cross-region-crypto-networks:hubs}.
\begin{tutsteps}
\item \textbf{Hubs.} The region cities are rows of \texttt{data/networks/sites.csv}; \texttt{nw\_race.hub\_rows()} computes geodesics, floors and
  \omterm{def:nw:cross-region-crypto-networks:infl}{path inflation} against the published round trips.
\item \textbf{Routes.} \texttt{ROUTES} holds three routes from published medians and assumed tails; \texttt{first\_arrival}
  (\cref{lst:nw:cross-region-crypto-networks:race}) draws them correlated.
\item \textbf{Gain and price.} \texttt{race\_rows()} and \texttt{break\_even()}.
\item \textbf{Duplicates.} \texttt{duplicates} applies the open-only rule; \texttt{duplicate\_rows()} sweeps the fill time.
\end{tutsteps}
\textbf{What to change next.} Give the backbone route to Frankfurt its published 224~ms and race it against the specialist; set the rule to
``ever'' and see the duplicates disappear.
\end{tutorial}

\section{Build: the route racer}\label{bld:nw:cross-region-crypto-networks:routerace}

\begin{build}
\bfield{Purpose} The latency, cost and risk of sending messages on several routes between regions: the input of the cross-region rows of chapter~29's
plan.
\bfield{Interface} \texttt{firm\_routerace}: \texttt{Route}, \texttt{first\_arrival}, \texttt{percentile\_gain}, \texttt{break\_even\_price},
\texttt{duplicates}, \texttt{path\_inflation}.
\bfield{Rules} Correlation is explicit; route draws use common random numbers across route counts; duplicate handling follows a stated venue rule;
simulations say so.
\bfield{Acceptance tests} \texttt{code/firm/routerace/tests/}: the marginals' median and 99th percentile, the first arrival, no gain at perfect
correlation, the duplicate rules by hand, \omterm{def:nw:cross-region-crypto-networks:infl}{path inflation} of a perfect route.
\bfield{Stretch} Fit route delays and their correlation from logged acknowledgements; a sender that races cancels and sends orders once.
\end{build}

\begin{omsources}
\item AWS for Industries blog, ``Ultra-low-latency cross-Region crypto trading with Avelacom and AWS'' (June 2026); AWS Regions documentation.
\item RFC 4271 (BGP-4); Binance spot API documentation (client order identifiers and duplicate orders).
\end{omsources}

\section{Exercises}

\begin{exercise}[$\star$]\label{exo:nw:cross-region-crypto-networks:1}
What is the \omterm{def:nw:cross-region-crypto-networks:infl}{path inflation} of the backbone and of the specialist from Tokyo to Frankfurt?
\end{exercise}

\begin{exercise}[$\star$]\label{exo:nw:cross-region-crypto-networks:2}
Why does BGP not choose the fastest path?
\end{exercise}

\begin{exercise}[$\star$]\label{exo:nw:cross-region-crypto-networks:3}
How much round trip does the specialist save over the backbone from Tokyo to each European hub?
\end{exercise}

\begin{exercise}[$\star\star$]\label{exo:nw:cross-region-crypto-networks:4}
Two routes share a subsea cable for half their length. What does that do to the value of racing them?
\end{exercise}

\begin{exercise}[$\star\star$]\label{exo:nw:cross-region-crypto-networks:5}
Which messages can be raced safely under an open-only duplicate rule, and which cannot?
\end{exercise}

\begin{exercise}[$\star\star$]\label{exo:nw:cross-region-crypto-networks:6}
A strategy trades Binance in Tokyo against a European venue. Where should its engine run?
\end{exercise}

\begin{exercise}[$\star\star\star$]\label{exo:nw:cross-region-crypto-networks:7}
\emph{Coding.} With \texttt{firm.routerace}, find the correlation above which the third route saves less than \qty{100}{\micro\second} at the 99th
percentile.
\end{exercise}

\begin{exercise}[$\star\star\star$]\label{exo:nw:cross-region-crypto-networks:8}
\emph{Find the flaw.} ``We race every order on three routes; the venue rejects duplicates, so it is free latency.''
\end{exercise}

\section{Problem: Three Routes to Singapore}

\begin{problem}[{Weekend problem --- what racing is worth}]\label{pb:nw:cross-region-crypto-networks:1}
A firm in Tokyo sends orders to a venue in Singapore on up to three routes: a specialised network (\qty{32.7}{\milli\second} one way, median), the
provider's backbone (\qty{34.0}{\milli\second}) and the internet (\qty{35}{\milli\second}); the two private routes' 99th percentiles are 15\% above
their medians, the internet's 30\%. A microsecond of 99th percentile is worth 10~USD a month to it.

\medskip\noindent\textbf{Part I --- The floors.}
\begin{enumerate}
\item What are the geodesic, the vacuum and the fibre round-trip floors between Tokyo and Singapore?
\item What is each published route's \omterm{def:nw:cross-region-crypto-networks:infl}{path inflation}?
\item How much of the specialist's round trip is above the fibre floor?
\item Why is the internet route's median an assumption rather than a published figure?
\end{enumerate}

\medskip\noindent\textbf{Part II --- The race.}
\begin{enumerate}[resume]
\item What does racing two routes save at the 99th percentile for $\rho = 0$, 0.5 and 0.9?
\item And three routes?
\item What is the third route worth a month at each correlation?
\item What does racing do to the median?
\end{enumerate}

\medskip\noindent\textbf{Part III --- The duplicates.}
\begin{enumerate}[resume]
\item Under an open-only rule, how many late copies of an order that fills on arrival are executed?
\item And of an order that takes a millisecond, four, eight to fill?
\item Which of the firm's messages would you race?
\item What rule would you ask the venue for?
\end{enumerate}

\medskip\noindent\textbf{Part IV --- The verdict.}
\begin{enumerate}[resume]
\item State the \emph{named result}: the 99th-percentile gain of racing two or three routes with correlation $\rho$, and the price per route above which
  racing does not pay.
\item Which input moves the answer most?
\item How would you estimate $\rho$?
\item What happens to $\rho$ when a cable is cut?
\item Would you race to Frankfurt, where the backbone is 88~ms slower than the specialist?
\item What does the firm's position inside each region (chapter~17) do to all of this?
\item Where does this go in the plan of chapter~29?
\item In one sentence: when does racing pay?
\end{enumerate}
\end{problem}

\section{Interview questions}

\begin{interviewq}[$\star$ \iqroles{developer}]\label{iq:nw:cross-region-crypto-networks:1}
Why is the round trip from Tokyo to Frankfurt over a \omterm{def:nw:cross-region-crypto-networks:bgp}{cloud backbone} more than twice the physical floor?
\end{interviewq}

\begin{interviewq}[$\star\star$ \iqroles{developer}]\label{iq:nw:cross-region-crypto-networks:2}
When would you pay for a specialised network between two cloud regions?
\end{interviewq}

\begin{interviewq}[$\star\star$ \iqroles{developer, researcher}]\label{iq:nw:cross-region-crypto-networks:3}
You send the same message on two routes. What does it buy you, and what does it depend on?
\end{interviewq}

\begin{interviewq}[$\star\star$ \iqroles{developer}]\label{iq:nw:cross-region-crypto-networks:4}
What is an \omterm{def:nw:cross-region-crypto-networks:idem}{idempotency key}, and how can it fail to protect you?
\end{interviewq}

\begin{interviewq}[$\star\star$ \iqroles{trader}]\label{iq:nw:cross-region-crypto-networks:5}
How does a hundred milliseconds between Tokyo and Europe shape a cross-venue crypto strategy?
\end{interviewq}

\begin{interviewq}[$\star\star\star$ \iqroles{developer, researcher}]\label{iq:nw:cross-region-crypto-networks:6}
Design an order-sending layer that races routes without ever executing an order twice.
\end{interviewq}
